% SPDX-FileCopyrightText: 2016-2026 Lari Natri <lari.natri@iki.fi>
% SPDX-License-Identifier: GPL-3.0-or-later
%
% Attached text/music note convenience commands and capo presentation.
% This file is part of the 'ulsbs' package.
%
% Load: ulsbs.sty, after song metadata and before repeats in the content phase.
% Requires: book visibility/transposition policy, xparse, expl3 and xcolor.
% Verse/chord/music owners later supply attachednote, transpose, music-comment
% guarding and verse separation before these command bodies are executed.
% State: owns convenience definitions and the capo override, not attachment
% machinery or display policy. Inherits @ as a letter with no enclosing group;
% retains the original expl3 span. Only definitions execute at load time.
% Public command arguments and verse-separation spacing resolve at use time;
% capo source transposition remains independent of display permission, and no
% new defaults or public override conventions are introduced by this module.

\ProvidesFile{ulsbs-notes.code.tex}[ULSBS attached notes and capo presentation]

% === Attached text and music notes ===

% \attachednote displays a note on a new line with attachment-aware spacing;
% the verse owner supplies its implementation. It is distinct from
% \textnote, the plain textual-note command.
%
% If given from within a stanza, only small vertical space is
% added before and after the note. Otherwise the vertical space given as
% optional parameters is used.
%
% The note is attached to the material after it (default), or to the matter
% after it (if first optional star is present). This means that column breaks do
% not happen between the note and the stuff it is attached to. Attachment is
% ignored if called from within a stanza.
%
% If the note is attached to the material after it, and called from outside of
% of stanzas, the \glueverses macro is called to remove the \vspace the
% beginning of the next verse in the current song would give. This might cause
% a problem if there is material between this note and the next stanza. In such
% a case one can attach the note to the material before, or add the second star
% to signify that \glueverses is never called from within this macro.
%
% Usage:
%
%   \attachednote*[vspace before][vspace after]{text}*
%
%   1st star:      This optional star signifies that the note is attached to the
%                  stuff before it. If it is omitted, the note is attached to
%                  the stuff coming after it.
%   vspace before: The amount of vspace added before the note. This is passed
%                  to \vspace, and can contain macros, for example:
%                  half the configured verse separation. Default: 0pt.
%   vspace after:  The amount of vspace added after the note. This is passed
%                  to \vspace, and can contain macros, for example:
%                  half the configured verse separation. Default: 0pt.
%   text:          The text content of the note.
%   2nd star:      If present, this star means that \glueverses is not called in
%                  any case.
%
% NOTE: If optional parameters are omitted, also the empty bracket pair must
%       be left out! So no "[]".
%
%
% Note: \NewDocumentCommand, IfValueTF etc. come from xparse package, and allow
%       nice handling of parameters.


% Creates a textual note attached to the material preceding it.
\newcommand{\textnoteforprev}[1]{%
  \attachednote*[0.618034\ulsbs@verse@separation][\ulsbs@verse@separation]{#1}%
}

% Creates a textual note attached to the material coming after it.
% Starred version omits glueverses. The starred version omits glueverses:
% use the star if before 1st verse.
\NewDocumentCommand{\textnotefornext}{m s}{%
  \IfBooleanTF{#2}{%
    \attachednote[\ulsbs@verse@separation][0.618034\ulsbs@verse@separation]{#1}*%
  }{%
    \attachednote[\ulsbs@verse@separation][0.618034\ulsbs@verse@separation]{#1}%
  }%
}%

% Creates a textual note in the middle of preceding and consequent material.
% It is attached to the material coming after it, but does not call \glueverses.
\newcommand{\textnotemiddle}[1]{%
    \attachednote[\ulsbs@verse@separation][\ulsbs@verse@separation]{#1}*%
}

% Creates a textual note in the middle of preceding and consequent material
% with smaller vertical space around it.
% It is attached to the material coming after it, but does not call \glueverses.
\newcommand{\textnotemiddletight}[1]{%
    \attachednote[0.618034\ulsbs@verse@separation][0.618034\ulsbs@verse@separation]{#1}*%
}

% Musical prose uses one independent lazy guard, including placement and spacing.
% Creates a chord-colored comment attached to the material preceding it.
\newcommand{\musicnoteforprev}[1]{%
  \ulsbs@musiccomment{\textnoteforprev{\color{ulChordColor}#1}}%
}

% Creates a chord-colored comment attached to the material coming after it.
% The trailing star omits glueverses: use it before the first verse.
\NewDocumentCommand{\musicnotefornext}{m s}{%
  \ulsbs@musiccomment{%
    \IfBooleanTF{#2}{%
      \textnotefornext{\color{ulChordColor}#1}*%
    }{%
      \textnotefornext{\color{ulChordColor}#1}%
    }%
  }%
}

% Creates a textual note in the middle of preceding and consequent material.
% It is attached to the material coming after it, but does not call \glueverses.
% The independent comment guard applies before the chord-colored note is placed.
\newcommand{\musicnotemiddle}[1]{%
  \ulsbs@musiccomment{\textnotemiddle{\color{ulChordColor}#1}}%
}

% Rewrite capo presentation to attach to the material coming after it, and to
% make it smaller than the default. Only shown if chords are enabled.
% Transposition is a source operation, never gated by display permission. The
% non-transposing instruction is chord-specific, independent of comment policy.
\ExplSyntaxOn
\DeclareDocumentCommand \capo {m}
  {\legacy_if:nTF {ulTransposeCapos}
    {\transpose{#1}}
    {\legacy_if:nT {ulShowChords}
      {\attachednote[\ulsbs@verse@separation][0.618034\ulsbs@verse@separation]
        {\color{ulChordColor}\scriptsize capo~#1}*}}}
\ExplSyntaxOff
